% This template was created by Chris Hoofnagle for 
% the Berkeley Center for Law & Technology. If you 
% work for UCB, you are welcome to use it.


\documentclass[letterpaper,12pt,openbib]{report}

% This is a "report" document class, which means that it has chapters
% and formatting like a book, but fundamentally is a design for
% one-sided printing. There are many options to change the underlying
% design---you can change the font, paper size, etc with these 
% options below. Just put the options in the braces of 
% \documentclass

% Documentclass: Report Options: a4paper, a5paper, b5paper, letterpaper, legalpaper, landscape, 11pt, 12pt, twoside, twocolumn, notitlepage, openright, draft, fleqn, leqno, openbib

%\usepackage[utf8]{inputenc}
\usepackage{graphicx}
\usepackage{lipsum} % generates lorem ipsum filler text
\usepackage[most,many,skins,breakable]{tcolorbox}
\usepackage{xcolor}
\usepackage{tikz}
\usepackage{tikzpagenodes}
\tcbuselibrary{skins,breakable,xparse}
\usepackage{fancyhdr}
\usepackage[default,scale=0.95]{opensans}
\usepackage[hidelinks]{hyperref} % the hidelinks option prevents those terrible boxes from appearing in your print pdf
\usepackage[backend=biber,language=american,style=verbose]{biblatex}
\addbibresource{bibliography.bib}
\usepackage[headheight=15pt,hmarginratio=1:1]{geometry}

%%%%%%%%%%%%%%%%%%%%%%%%
%
% Berkeley Colors - Global (changing these will alter the whole
% document)
%
%%%%%%%%%%%%%%%%%%%%%%%%

% Deep blue and gold
\definecolor{berkeleyblue}{HTML}{003262}
\definecolor{berkeleygold}{HTML}{FDB515}
% Light blue and gold
\definecolor{founderblue}{HTML}{3B7EA1}
\definecolor{founderyellow}{HTML}{C4820E}

%%%%%%%%%%%%%%%%%%%%%%%%
%
% Headings
%
%%%%%%%%%%%%%%%%%%%%%%%%

\fancyhf{}
% If you'd like, you can put the title of the report in the header
% Just erase title of report if not

\small\fancyhead[R] {\color{founderblue}Title of Report ~ \thepage}
% This adds the line below the header. You can make it thicker or thinner by editing the 0.4 value
\renewcommand{\headrulewidth}{0.4pt}
% This template only uses a header because a document with footnotes
% looks strange with a footer. But if you want a footer, just 
% uncomment the lines below
%\small\fancyfoot[R] {\color{founderblue}Title of Report ~ \thepage}
%\renewcommand{\footrulewidth}{0.0pt}

%%%%%%%%%%%%%%%%%%%%%%%%
%
% Customizing Fonts & Title Formats - Global
%
%%%%%%%%%%%%%%%%%%%%%%%%

% This section specifies the fonts and their sizes and 
% colors for the headings. Notice that the package gives
% you serif headings and san serif body text.

\usepackage{titlesec}
\usepackage{fontspec}

% The headings are Berkeley's Old Style Fonts!
% I don't like the new fonts, which are licensed from
% Adobe --- why would you do that?

\newfontfamily\headingfont{UCBerkeleyOS.otf}[
BoldFont = UCBerkeleyOSBold.otf,
ItalicFont = UCBerkeleyOSItalic.otf,
BoldItalicFont = UCBerkeleyOSBoldItalic.otf]

% This controls the chapter headings

\titleformat{\chapter}{\fontsize{30}{60}\headingfont}{\color{founderblue}\thechapter\fontsize{90}{60}}{10pt}{\thispagestyle{empty}\color{founderblue}\Huge\bfseries}

% This controls the section and subsection headings in the chapters
% Notice that the headings in the chapter are numbered, e.g.
% 1.1 A Lesson in Latin. If you don't want section numbering,
% remove \thesection and/or \thesubsection

\titleformat{\section}{\Large\headingfont}{}{0pt}{\color{founderblue}{\thesection} \ }[{\titlerule[0.8pt]}]

\titleformat{\subsection}{\large\headingfont}{}{0pt}{\color{founderblue}{\thesubsection} \ }[{\titlerule[0.8pt]}]

% This customizes the sections in the executive summary so that they
% will not be numbered.

\titleformat{name=\section,numberless}
{\Large\headingfont}{}{-4pt}{\color{founderblue} \ }[{\titlerule[0.8pt]}]

\titleformat{name=\subsection,numberless}
{\large\headingfont}{}{-4pt}{\color{founderblue} \ }[{\titlerule[0.8pt]}]

%%%%%%%%%%%%%%%%%%%%%%%%
%
% Sidebar Environment: mdframed
%
%%%%%%%%%%%%%%%%%%%%%%%

\usepackage{amstext}
\usepackage{mathtools}
\usepackage[framemethod=TikZ]{mdframed}


\mdfdefinestyle{custom}{frametitle={\colorbox{berkeleygold}{\space You must specify a title for your box \space}},
     frametitlefont=\textbf,
     backgroundcolor=founderblue!20,
     linecolor=founderblue!70,
     innertopmargin=10pt,
     frametitleaboveskip=-\ht\strutbox,
     skipabove=\topskip,
     skipbelow=\topskip}




%\newenvironment{sidebar}[1]
% \mdfsetup{%
% % Replace berkeleygold with white if you'd like
%     frametitle={\colorbox{berkeleygold}{\space#1\space}},
%     frametitlefont=\textbf,
%     backgroundcolor=founderblue!20,
%     linecolor=founderblue!70,
%     innertopmargin=10pt,
%     frametitleaboveskip=-\ht\strutbox,
% %    frametitlealignment=\left,
%     skipabove=\topskip,
%     skipbelow=\topskip
%     }%
%   \begin{mdframed}%
%   }
%   {\end{mdframed}}

%%%%%%%%%%%%%%%%%%%%%%%%
%
% Sidebars with tcolorbox -- Mybox
%
%%%%%%%%%%%%%%%%%%%%%%%

\newtcolorbox{mybox}[2][]{enhanced, interior hidden, 
coltitle=black, fonttitle=\bfseries\headingfont\large,
attach boxed title to top left={yshift=-2.5mm},
boxed title style={empty, size=small, top=1mm, bottom=2pt},
%boxed title=0.5\linewidth,
frame code={
\path (title.east|-frame.north) coordinate (aux);
\path[draw=founderblue, fill=founderblue!20, line width=0.5mm, rounded corners]
(frame.west) |- ([xshift=-2.5mm]title.north east) to[out=0, in=180] ([xshift=7.5mm]aux)-|(frame.east)|-(frame.south)-|cycle;  
},
title={#2},#1}

%%%%%%%%%%%%%%%%%%%%%%%%
%
% General Guidance
%
%%%%%%%%%%%%%%%%%%%%%%%

%\flushbottom

%The \flushbottom declaration makes all text pages the same height, adding extra vertical space when necessary to fill out the page.

%\onecolumn
%The \onecolumn declaration starts a new page and produces single-column output.

%\raggedbottom
%The \raggedbottom declaration makes all pages the height of the text on that page. No extra vertical space is added.

%\twocolumn
%The \twocolumn declaration starts a new page and produces two-column output.

%%%%%%%%%%%%%%%%%%%%%%%%
%
% Document Content Starts Here
%
%%%%%%%%%%%%%%%%%%%%%%%


\title{Put title here, and in the title.tex file}

\begin{document}

% The title page appearance is controlled by a separate file named titlepage.tex
\begin{titlepage}
\pagecolor{founderblue}
       \vspace*{1cm}


% This controls the title
% You can reposition this element by editing the numeric
% values in \node (1,1). The first value controls the 
% horizontal placement, the second the vertical.

\tikz[overlay, remember picture] \node at (1,1) {
\begin{tcolorbox}[colback=berkeleygold,width=15cm,halign=right,boxrule=0pt]%%
\headingfont{\fontsize{40}{60}\selectfont This is my title:}

\end{tcolorbox}
};

% This controls the subtitle
\tikz[overlay, remember picture] \node at (1,-1.4) {
\begin{tcolorbox}[colback=berkeleygold,width=15cm,halign=right,boxrule=0pt]%%
\headingfont{\fontsize{30}{60}\selectfont This is my subtitle}

\end{tcolorbox}
};



   

\vspace*{1cm}
% This controls the author block. You can add or reduce the
% number of authors by simply deleting text and the \\. Each
% \\ gives you a new line
\tikz[overlay, remember picture] \node at (-4,-8) {
\begin{tcolorbox}[colback=berkeleygold,width=15cm,halign=right,boxrule=0pt]%%
  \headingfont\Large Author 1 \\ Author 2 \\ Author 3
\end{tcolorbox}
};
  
% This controls the date block. \today produces today's date.
% You can erase \today and replace it with something else
\tikz[overlay, remember picture] \node at (-4,-12) {
\begin{tcolorbox}[colback=berkeleygold,width=15cm,halign=right,boxrule=0pt]%%
  \headingfont\Large\today
\end{tcolorbox}


};
    
% This places the front cover logo. Replace logo.png with
% your logo
% You can reposition the logo by changing the node values
% 6.5 places the image in the center. The second number
% controls its vertical placement
\tikz[overlay, remember picture] \node at (6.5,-15) {
\includegraphics[width=0.7\textwidth]{logo.png}  


};


\end{titlepage}

% This makes the report have a white background
\pagecolor{white}

% Just comment out the next lines if you 
% do not want a table of contents

% This adds the Executive Summary to the TOC w/o a number
\addcontentsline{toc}{chapter}{Executive Summary}
\tableofcontents
% Note: to have a conclusion in the TOC without a number
% the addcontentsline command appears after \include conclusion
% below

% The content of the report is in subfiles that are
% imported below with the include command
% You can create more "chapters" by copying the chapter.tex
% files and using \include{chapterX} to include them in
% the report

\chapter*{Executive Summary}
\pagestyle{fancy}

% The Asterix after chapter here excludes executive summary from the table of content and ensures it doesn't get a number (e.g. chapter 1)



\section*{Section title for exec summary}
% The Asterix after section causes it to not be included in the table of contents

\lipsum[1]

\subsection*{Sub section title for exec summary}

\lipsum[2-3]

\chapter{Welcome to chapter one!}
\pagestyle{fancy}

You can do normal footnotes\footnote{Hello, here is a footnote.}

Or you can manage your footnotes through a bibliography file.\footnote{\cite{warren1890right}}

\bigskip

\lipsum[4]

\bigskip

\begin{mybox}{Hello World!}

This is the mybox environment. You might like this for definitions.

However, the problem with mybox is that they cannot break across pages :( So, they have to be pretty short. Otherwise use the colorbox

\end{mybox}

\bigskip

\section{A lesson in Latin!}

\lipsum[5-8]



\chapter{Welcome to chapter two}
\pagestyle{fancy}

\lipsum[9]



\begin{mdframed}[style=custom,frametitle=\colorbox{berkeleygold}{\space Put your title here \space}]

You might like to use these boxes for essays or to break up text. Unlike the mybox environment, mdframed will break across pages.

Don't like the colors?

You can change them by editing the sidebar section of report.tex. In mdfdefinestyle, look for \textbackslash colorbox\{berkeleygold\} to alter the title background. If you want no color, just change berkeleygold to white. If you want to dilute the color a bit, you could add a !30 to the end of berkeleygold, as in berkeleygold!30

The background color is set by backgroundcolor and it is currently founderblue!20. You can deepen it by increasing the number or lighten it by reducing the number to 10 or even 5.

\end{mdframed}


\begin{mdframed}[style=custom,frametitle=\colorbox{berkeleygold}{\space This is my second box \space}]

Please note, you have to specify your title in each box in the frametitle option. Just put your title between the \textbackslash space commands.

\end{mdframed}


\chapter{Welcome to chapter three}
\pagestyle{fancy}

\begin{tcolorbox}[breakable,title={An example colorbox},
colback=founderblue!5!white,
colframe=founderblue!75!black,
fonttitle=\headingfont\bfseries\large]
Your text goes here. The colors are based on Berkeley's Founder Blue.
\tcbsubtitle[before skip=\baselineskip]%
{You can also have subboxes}
Don't like the colors? You can change them here (replace founderblue with the color you like), or you can change them globally in the colors section of report.tex
\tcbsubtitle[before skip=\baselineskip]%
{Multiple rows}
If you like it.
\end{tcolorbox}







%\include{chapter4}
%include{chapter5}
\chapter*{Conclusion}
\pagestyle{fancy}

We're using chapter 4 as the conclusion. If you want to have more chapters, just create more files with chapter[number].tex. Make the last one the conclusion. In order to omit the chapter number, just place an asterix before the chapter name like so:

\textbackslash chapter*\{Conclusion\}

% This adds the conclusion to the TOC without a number
\addcontentsline{toc}{chapter}{Conclusion}
% It is possible to add a bibliography here
\printbibliography
% This is the back cover
% The page background is Berkeley's dark blue color


\pagecolor{berkeleyblue}
\pagestyle{empty}

% This places the back cover logo. Replace logo.png with
% your logo
\tikz[overlay, remember picture] \node at (6.5,-14) {
\includegraphics[width=0.7\textwidth]{logo.png}    
};

% This places those lines of text. Just erase Text 1, etc if you
% don't want any text

\tikz[overlay, remember picture] \node at (6.5,-17.7) {
\begin{tcolorbox}[colupper=white,colback=berkeleyblue,width=10cm,halign=left,boxrule=0pt]%%
  \headingfont\Large 
  Text 1 \\ 
  Text 2 \\ 
  Text 3
\end{tcolorbox}
};



% This code creates decorative `colorboxes' for text
% illustrations and the like. Just uncomment it and place
% it in a chapter

% \begin{tcolorbox}[breakable,title={An example colorbox},
% colback=founderblue!5!white,
% colframe=founderblue!75!black,
% fonttitle=\headingfont\bfseries\large]
% Your text goes here.
% \tcbsubtitle[before skip=\baselineskip]%
% {You can also have subboxes}
% with extra content
% \tcbsubtitle[before skip=\baselineskip]%
% {Multiple rows}
% If you like it.
% \end{tcolorbox}


% This code creates a decorative text box. 
% Just uncomment it and place
% it in a chapter

% \begin{mybox}{Hello World!}

% This is the mybox environment.

% However, the problem with mybox is that they cannot break across pages :( So, they have to be pretty short. Otherwise use the colorbox

% \end{mybox}





\end{document}

